\section{真机与模拟机对比}
\label{sec:comparison}

图~\ref{fig:D08D09}(c,d) 比较模拟机 $p=1,2$ 与真机 $p=1$ 的结果。
百花处理器上的测量以重复采集单元组织。
每个单元在 $\delta\in\{0.85,-2\}$ 两条扫描线上各取 $15$ 个 $s$ 值，共 $30$ 点，每点 $1024$ shots，并额外执行一对读出标定电路。
本文分析一轮 $500$ 个单元。

真机观测量来自同一批读出缓解分布。
每个候选窗口以其 pooled 的 $30$ 点均值拟合两线共用最大值，定义
\begin{equation}
\hat S_{\rm hw}=\frac{S_{\rm hw}}{S_{\max}},\qquad
\hat O_{\rm hw}=\frac{-O_{{\rm str},{\rm hw}}}{O_{\max}}.
\label{eq:hardware-scales}
\end{equation}
所选窗口的 $S_{\max}=4.57$、$O_{\max}=0.089$。
模拟机保持式~\eqref{eq:fixed-scales} 的固定尺度。
分别将两组标记代入式~\eqref{eq:Q} 的相同系数，组合后的 $Q$ 不再归一化。
真机尺度将两类关联量置于共同尺度上，用于比较其参数依赖。

比较指标为 $30$ 点平均绝对 $Q$ 距离 $d_{\rm raw}$ 与 Pearson 相关系数 $r$。
$p=1$ 模拟曲线用于选窗，$p=2$ 仅作图示对照。
候选窗口为 $500$ 单元序列中长度至少为 $10$ 的全部连续偶数长度子段，各窗口重新拟合最大值。
选择 $r$ 最大者，并列时依次取 $d_{\rm raw}$ 更小、更长、更早的窗口。
最终窗口为 $[222,238)$，含 $16$ 个单元，$d_{\rm raw}=0.30$、$r=0.86$。
误差棒为该窗口内逐单元 $Q$ 值的样本标准差除以 $\sqrt{16}$，所选尺度保持固定。

$\delta=-2$ 线上，测得的 $Q$ 从小 $s$ 区间的正值转为大 $s$ 区间的负值。
$\delta=0.85$ 线上，小 $s$ 区间为正值，中间区间出现近零和负值。

为保留振幅和符号信息，另检查除以拟合 $O_{\max}$ 之前的 pooled string 量。
$s=0.99$、$\delta=0.85$ 时，固定尺度的 $\hat O_{\rm str}=-O_{\rm str}$ 在全部 $500$ 单元中为 $0.043\pm0.005$（单元间标准误），对应模拟值为 $+0.95$。
